% ============================================================
% P05 — Verifikasi Email
% ============================================================
\flowmeta{P05}{Verifikasi Email}{Pelanggan}

\begin{flowsection}{Tujuan}
Memverifikasi alamat email akun agar status akun menjadi terverifikasi dan
pelanggan dapat menerima notifikasi penting.
\end{flowsection}

\begin{flowsection}{Prasyarat}
Sudah mendaftar (P01) dan menerima email verifikasi.
\end{flowsection}

\begin{flowsection}{Langkah Verifikasi}
\begin{enumerate}
  \item Buka email verifikasi, klik tautan verifikasi.
  \item Sistem membuka halaman \texttt{/email-verified} dan menandai email
        sebagai terverifikasi.
  \item Bila belum menerima email, buka beranda/akun saat masuk $\rightarrow$
        klik \textbf{Kirim Ulang Verifikasi} pada banner.
\end{enumerate}
\end{flowsection}

\shot{gambar/pelanggan/p05-1.png}{Halaman hasil verifikasi email}{fig:p05-1}

\begin{flowsection}{Hasil yang Diharapkan}
Email ditandai terverifikasi; banner \textbf{Email Belum Diverifikasi} hilang dan
pengguna dapat menerima notifikasi pesanan.
\end{flowsection}

\shot{gambar/pelanggan/p05-2.png}{Banner ``Email Belum Diverifikasi'' saat masuk dengan akun yang belum terverifikasi}{fig:p05-2}

\begin{flowsection}{Penanganan Error dan Kasus Khusus}
\begin{itemize}
  \item Tautan kedaluwarsa/tidak valid $\rightarrow$ minta kirim ulang.
  \item Email verifikasi tidak sampai (mis. email masuk spam) $\rightarrow$
        gunakan tombol kirim ulang.
  \item Pengiriman email gagal di sisi server $\rightarrow$ dicatat di log;
        pengguna tetap dapat mencoba kirim ulang.
\end{itemize}
\end{flowsection}

\begin{flowsection}{Kaitan Modul}
FE \texttt{EmailVerifiedPage.jsx}, banner verifikasi di \texttt{App.jsx}. BE
\texttt{AuthController@verifyEmail}, \texttt{@resendVerificationEmail}. Rute
\texttt{GET /api/auth/email/verify/\{id\}/\{hash\}}, \texttt{POST /api/auth/email/resend}.
\end{flowsection}

\docver{v1.0}{Sep 2026}
